\subsection{The rank-two constituents and the remaining rigidity}
\label{sec:support-type-two-rigidity}

For the original good datum, the remaining constituents are covered by
Kato's theorem at a noncentral critical twist.  The fact that the
classical weight is odd is decisive: no central-value hypothesis is
required.

Retain the original datum $D$, its imaginary quadratic field $K$, odd
prime $\ell$, coefficient ring $O$, character
$\Psi=\varepsilon^e\chi$, and all its original good-prime hypotheses.
In particular $e\ge1$, $\chi$ has order prime to $\ell$,
$K$ is unramified at $\ell$, and $\ell\nmid h_K$, where $h_K$
is the class number.  These are original hypotheses, not consequences
of having finitely many target characters.  The full original
good-data hypotheses, including their conductor coprimality, are retained.
No assertion that arbitrary first-family data satisfy these hypotheses
is made: that family's original end condition requests only finite-end
control and rank-three automorphy.  The rank-two end condition and the
two-sided rigidity consumer quantify precisely these good data
\cite{support-type-two-source}.

Let $f$ be a local characteristic-zero point of any original global
ring $S$, with its original selected local ring.  Write $r=r_f$ and
suppose it is type II.  Thus on $G_K$ its matrix part has the form
\[
 r(g)_{\rm mat}=
 \begin{pmatrix}a(g)&0&b(g)\\0&d(g)&0\\c(g)&0&d'(g)\end{pmatrix},
 \qquad
 A(g)=\begin{pmatrix}a(g)&b(g)\\c(g)&d'(g)\end{pmatrix},          \tag{RR1}
\]
and the two-dimensional representation $A$ is absolutely irreducible.
Use one finite coefficient field $E$ and its invariant original
lattice, supplied by the finite-neighbourhood argument.

\begin{lemma}[The original rank-two representation]
\label{lem:support-type-two-representation}
In (RR1), $d=1$ and $\det A=1$ on $G_K$.  There is a continuous
representation $W:G_{\mathbf Q}\to\mathrm{GL}_2(E)$ extending $A$,
with
\[
 W(c)=C=\operatorname{diag}(-1,1),\qquad \det W=\eta.            \tag{RR2}
\]
It is absolutely irreducible on every open subgroup and on the kernel
of the cyclotomic character.  It is the original rank-two extension
required at a type-II point.
\end{lemma}

\begin{proof}
Both $d$ and the determinant of the full matrix representation on
$G_K$ are characters congruent to one.  At $\ell$, the ordered
ordinary-flag lemma~\ref{lem:support-integral-ordinary-flag} supplies
the flag over the original coefficient ring.  Its Hensel spectral
basis at the decomposition-group mark respects the outer block and
the middle line in (RR1): the middle eigenvalue is the unique root
congruent to one.  Hence the middle inertia character is exactly
one, the outer two are $\Psi,\Psi^{-1}$, and the full determinant
is one on inertia.

At any other finite prime the original equations kill inertia in
$\ker\eta\cap\ker\Psi$.  Since $\varepsilon$ is unramified there,
the remaining quotient of inertia in $G_K$ has order dividing
$\operatorname{ord}(\chi)$, prime to $\ell$.  The characters $d$
and the full determinant therefore have prime-to-$\ell$ order there
and are congruent to one, so are one.  They are consequently everywhere
unramified characters of $G_K$.  Class field theory makes their orders
divide $h_K$.  Since $\ell\nmid h_K$, reduction to one makes both
characters trivial.  This proves $d=\det A=1$.

Put $J_2=\left(\begin{smallmatrix}0&1\\1&0\end{smallmatrix}\right)$.
The original polarization and the normalization at $c$ give
\[
 A(cgc^{-1})=J_2A(g)^{-\mathsf t}J_2=CA(g)C,\qquad g\in G_K.
                                                               \tag{RR3}
\]
The second equality uses $\det A=1$.  Define $W(g)=A(g)$ on
$G_K$ and $W(gc)=A(g)C$ there.  Equation (RR3) and $c^2=1$ prove
multiplicativity, and give (RR2).  All its matrices and inverses
remain in the same coefficient lattice.

The proof of Lemma~\ref{lem:support-type-one-strong-irreducibility}
has the following rank-two form.  Reducibility on an open normal
subgroup gives a scalar action or two permuted lines.  An ordinary
inertia element with $\Psi$ of infinite order excludes the scalar
case.  The local ordinary flag makes the local decomposition group
fix both lines.  A global element interchanging them has trace zero
after multiplication by $1$ and by the decomposition mark $x$.
Reducing these identities gives
\[
        \alpha+\alpha^{-1}=0,\qquad
        \alpha\beta+(\alpha\beta)^{-1}=0,
        \quad\beta=\psi(x).
\]
They imply $\alpha^2=-1$ and $\alpha^2\beta^2=-1$, contradicting
$\beta^2\ne1$.  Thus $A$ is irreducible on every open subgroup of
$G_K$, and $W$ is irreducible on every open subgroup of
$G_{\mathbf Q}$ by intersection with $G_K$.
The connected reductive monodromy and derived-group argument in
that lemma then proves irreducibility on every kernel with abelian
quotient, in particular the cyclotomic kernel and the cyclotomic
kernel after any finite extension.
\end{proof}

\begin{lemma}[The actual odd-weight newform]
\label{lem:support-type-two-newform}
There is a classical normalized cuspidal newform $g$ of weight
$k=2e+1$ whose representation in the original geometric-Frobenius
normalization is
\[
                       V_g=W\varepsilon^{-e}.                  \tag{RR4}
\]
Consequently $W=V_g(e)$, with exactly Kato's convention for this twist.
\end{lemma}

\begin{proof}
Use the already constructed original auxiliary frame and its abelian
CM field $L/\mathbf Q$.  The normalized-lift transfer gives a
rank-three normalized representation.  Taking its outer block gives
the original rank-two normalized representation over the same $O_E$.
Indeed the middle line is trivial by the preceding lemma; the
polarization restricts to $J_2$, the marked diagonal relation
restricts to its two outer entries, and the residual characters are
$\psi,\psi^{-1}$.  The local Hensel basis in the preceding proof
restricts to an integral basis of the corner, congruent to the
identity, with ordered inertia characters $\Psi,\Psi^{-1}$.
Their two ordered factors lower this two-step flag successively,
giving the original rank-two local product equations.  Continuity
and the quotient identities are inherited without changing the
coefficient ring.

Use the actual algebraic CM character $\kappa_2$ from
Theorem~\ref{thm:actual-cm-twist}.  It has the original polarization
identity $\kappa_2\kappa_2^c=\varepsilon^{-1}$ on $G_L$.
The irreducibility just proved permits the rank-two part of
Theorem~\ref{thm:support-finite-points} to be applied.  The resulting
point of the actual Hecke quotient has determinant that of
$W|_{G_L}\kappa_2$.  If the labeled Hodge--Tate weights of
$\kappa_2$ are $a_\tau$, its two ordered labels are
$a_\tau-e,a_\tau+e$, so are strictly increasing.  The point is
arithmetic.  The full classicality in
Theorem~\ref{thm:selmer-ordinary-families} makes
$W|_{G_L}\kappa_2$ regular algebraic, conjugate self-dual and
cuspidal automorphic.

Algebraic twisting preserves polarized automorphy by
BLGGT, Lemma~2.2.1; apply it to $\kappa_2^{-1}$.
The multiplier after untwisting is trivial on $G_L$ and is $-1$
at complex conjugation: the transfer character of $\kappa_2^{-1}$
has value one there, because $c^2=1$.  Thus it is
$\eta|_{G_{L^+}}$.
The global representation $W$ is polarized by its alternating
determinant pairing with odd multiplier $\eta$.
The extension $L/\mathbf Q$ is finite, Galois and abelian, hence
solvable, and $W|_{G_L}$ is irreducible.  BLGGT, Lemma~2.2.2,
therefore descends automorphy to $W$ over $\mathbf Q$
\cite{support-type-two-blggt}.  These lemmas have no residual
irreducibility or image-size hypothesis.

The local flag makes the Hodge--Tate weights of $W$ equal to
$-e,e$, and $W(c)$ has eigenvalues $-1,1$.  After twisting by
$\varepsilon^{-e}$ the labels are $0,2e$.  The classical
description of regular algebraic cuspidal representations of
$\mathrm{GL}_2/\mathbf Q$ gives a holomorphic cuspidal newform of
weight $2e+1$ with this representation.  More explicitly the difference
of its two infinitesimal-character labels is $2e$, so its discrete
series weight is $2e+1$; the algebraic norm twist is fixed by the
labels $0,2e$.  Passing to the new vector fixes the normalized
newform and does not change the representation.  Its nebentype is
$\eta$, since its determinant is $\eta\varepsilon^{-2e}$.

There is no dual or inverse twist hidden in (RR4).  In the source,
at every good prime $q$, the characteristic polynomial at inverse
arithmetic Frobenius is
$X^2-a_q(g)X+\eta(q)q^{k-1}$.  Kato's formula (14.10.3), on
p.~240, uses the same inverse arithmetic Frobenius and the same
polynomial.  His determinant formula (14.10.2) is
$\varepsilon^{1-k}\eta^{-1}$, also the determinant just obtained.
Chebotarev and semisimple character theory therefore identify his
$V_g$ with (RR4), after a finite coefficient extension
\cite{support-type-two-kato}.

The original rank-two end condition permits any extension of $A$,
not just (RR2).  If $W'$ is another extension, Schur's lemma on
the absolutely irreducible restriction to $G_K$ gives
$W'(c)=\lambda W(c)$; the relation $c^2=1$ gives
$\lambda^2=1$.  Thus $W'=W$ or $W\eta$.  In the second case
the classical twist $g\otimes\eta$ has representation
$V_g\eta=W'\varepsilon^{-e}$.  It has the same weight and
is again cuspidal, so the full original extension quantifier is covered.
\end{proof}

\begin{lemma}[The two rank-two Selmer groups]
\label{lem:support-type-two-kato}
For the original $W$,
\[
          H_f^1(\mathbf Q,W)=H_f^1(\mathbf Q,W\eta)=0.            \tag{RR5}
\]
\end{lemma}

\begin{proof}
Apply Kato's Theorem~14.2(1) over the original finite abelian
extension $K/\mathbf Q$, with $r=e$ and $k=2e+1$.
Its inequalities are
\[
             1\le e\le2e=k-1,\qquad e\ne k/2.
\]
Thus for every coefficient place and every stable lattice $T$ in
$V_g(e)$, Kato's discrete Selmer group $\operatorname{Sel}(K,T)$
is finite.  His \S14.1 defines its local conditions as the images
of the same unramified and Bloch--Kato finite conditions used here.
The theorem includes forms with complex multiplication, treated
in his \S15, and excludes no coefficient prime.

Here is the precise passage from this finite discrete group to
rational vanishing.  Work first over a coefficient field of $g$ and
write $V=V_g(e)$, with stable lattice $T$.  Choose
$u\in I\cap G_K$ with $\Psi(u)$ of infinite order.
The matrix $B=V(u)-1$ is invertible, since its eigenvalues before
subtracting one are $\Psi(u),\Psi(u)^{-1}$.
For a cocycle $z$ representing a class in $H_f^1(K,V)$,
subtract the coboundary with primitive $B^{-1}z(u)$, making
$z(u)=0$.  Choose $a$ such that $\ell^a\operatorname{Sel}(K,T)=0$.
For every $n\ge a$, the class of $\ell^{-n}z$ modulo $T$ belongs
to this discrete Selmer group.  Consequently there is $t_n\in V/T$
with
\[
   \ell^{a-n}z(s)=(V(s)-1)t_n\pmod T,\qquad s\in G_K.
\]
At $s=u$ we have $Bt_n=0$ modulo $T$, so $t_n$ has a lift in
$B^{-1}T$.  Choose $b$ with $B^{-1}T\subseteq\ell^{-b}T$.
Because $T$ is stable, this gives
$z(s)\in\ell^{n-a-b}T$ for every $n\ge a$.
Their intersection is zero, so the normalized cocycle vanishes.
No compatibility among the $t_n$ is required.

This proves $H_f^1(K,V_g(e))=0$.  Finite coefficient extension
commutes with these cohomology groups and local kernels, as is
seen by choosing a basis of the coefficient extension; (RR4)
therefore gives $H_f^1(K,W)=0$.  Both $W$ and $W\eta$ restrict
to $W$ over $K$.  Characteristic-zero restriction injectivity
now gives (RR5).
\end{proof}

\begin{theorem}[The original type-II rigidity]
\label{thm:support-type-two-rigidity}
For every original good datum and every type-II local
characteristic-zero point of any of its original global rings,
\[
                    H_f^1(\mathbf Q,\operatorname{ad}r)=0.
                                                               \tag{RR6}
\]
Every original dual-number lift at that point is a coboundary.
Together with Theorem~\ref{thm:support-type-one-rigidity}, this
supplies the full two-sided characteristic-zero rigidity requirement
of the original good-data deformation argument.
\end{theorem}

\begin{proof}
We give the full coefficient-module decomposition.  The outer block
of the polarized adjoint has action $Y\mapsto A(s)YA(s)^{-1}$ on
$G_K$ and $Y\mapsto-J_2Y^{\mathsf t}J_2$ at $c$.  On traceless
matrices the latter is $Y\mapsto CYC$; on scalars it is minus one.
Thus this four-dimensional block is
$\operatorname{ad}^0W\oplus E(\eta)$.  The independent middle
diagonal line is another $E(\eta)$.

Write the remaining four entries as an outer column $u$ and an
outer row $v$.  Put
$\Omega=\left(\begin{smallmatrix}0&1\\-1&0\end{smallmatrix}\right)$
and $w=\Omega v^{\mathsf t}$.  On $G_K$, both $u$ and $w$ transform
by $A$.  At $c$, negative antitranspose gives
$(u,w)\mapsto(Cw,Cu)$.  Hence $u+w$ and $u-w$ give respectively
$W$ and $W\eta$; division by two recovers $u,w$.
The exact decomposition is consequently
\[
 \operatorname{ad}r\simeq
   \operatorname{ad}^0W\oplus E(\eta)^{\oplus2}
                  \oplus W\oplus W\eta.                        \tag{RR7}
\]
These are equivariant direct summands, so the finite local conditions
and Selmer groups decompose in the same way.

To treat the adjoint block, use the already obtained cuspidal
representation $W|_{G_L}\kappa_2$.  The rank-two irreducibility
lemma gives its irreducibility over $L(\mu_{\ell^\infty})$.
Its polarization matrix is $J_2$ and its multiplier on $G_{L^+}$
is $\varepsilon^{-1}$: its restriction to $G_L$ follows from
$\kappa_2\kappa_2^c=\varepsilon^{-1}$, and the value at $c$ is
$-1$.  Thorne's Theorem~6.1 applies in rank two.  Its adjoint is
exactly the restriction of the outer four-dimensional block above,
since scalar twisting does not change this action.  Restriction
injectivity descends its Selmer vanishing from $L^+$ to $\mathbf Q$
\cite{support-type-one-thorne}.

Since $E(\eta)$ is an equivariant direct summand of the outer
block, that same vanishing also kills the second copy of $E(\eta)$.
The two rank-two groups vanish by (RR5).  Equation (RR7) now proves
(RR6).

Finally every actual dual-number tangent satisfies the finite local
conditions by the same crystalline-component and finite-inertia
arguments used in Theorem~\ref{thm:support-type-one-rigidity}.
Those arguments are independent of irreducibility.  Equation (RR6)
makes its cocycle a coboundary, with all original dual-number and
local-ring quantifiers unchanged.
\end{proof}

This proves rational rigidity.  Neither Kato's qualitative
finiteness nor characteristic-zero corestriction bounds the torsion
in the original integral deformation modules uniformly as the
characters, levels and auxiliary fields vary.  That integral
estimate remains an independent frontier.
